\documentclass[12pt, a4paper]{article}
\usepackage[T1]{fontenc}
\usepackage{lmodern}
\usepackage{amsmath, amsthm, amssymb, mathtools}
\usepackage[margin=1in]{geometry}
\usepackage{xr-hyper}
\usepackage[colorlinks=true,linkcolor=blue,citecolor=blue]{hyperref}
\usepackage{cleveref}
\usepackage{graphicx, subcaption}
\usepackage{natbib}
\usepackage{booktabs}
\usepackage{enumerate}

% Code references: scaled monospace with line-breaking
\newcommand{\code}[1]{\texttt{\small #1}}

\newtheorem{theorem}{Theorem}[section]
\newtheorem{lemma}[theorem]{Lemma}
\newtheorem{proposition}[theorem]{Proposition}
\newtheorem{corollary}[theorem]{Corollary}
\newtheorem{definition}{Definition}[section]
\newtheorem{remark}{Remark}[section]
\newtheorem{conjecture}[theorem]{Conjecture}
\newtheorem{derivation}[theorem]{Derivation}
\newtheorem{observation}[theorem]{Observation}
\numberwithin{equation}{section}

\newcommand{\dd}{\mathrm{d}}
\newcommand{\seriestitle}{Why Rotating Fluids Make Polygons}


\externaldocument[I-]{../paper-1-mathematics/main}
\externaldocument[II-]{../paper-2-physics/main}
\externaldocument[III-]{../paper-3-gravity/main}
\externaldocument[IV-]{../paper-4-field-theory/main}
\externaldocument[V-]{../paper-5-s3-framework/main}
\externaldocument[VI-]{../paper-5-cosmology/main}
\externaldocument[VII-]{../paper-6-discussion/main}

\title{\seriestitle\\[6pt]\large Appendices: Technical Derivations and Open Problems}
\author{Gordon Speagle}
\date{}

\begin{document}
\maketitle

\begin{abstract}
This appendix collects the derivations of $C_1(\mathbf{H}^2, \xi)$ and $C_1(\mathbf{S}^2, \xi)$, code availability information, a derivation index for the series, and a list of open problems identified during this research.
\end{abstract}


\section{Derivation of $C_1(\mathbf{H}^2,\xi)$ and the near-$K=0$ expansion}
\label{app:c1_h2}

We work in the Poincar\'e disk of radius~$a$ and set $a = 1$ without
loss of generality (the formula is dimensionless in $\xi = r_E^2/a^2$).

\paragraph{H² Hamiltonian and moment map.}
In Poincar\'e disk coordinates the $N$-vortex Hamiltonian is
\citep{Kimura1999}
\begin{equation}\label{eq:H_hyp}
  H_{\mathrm{hyp}} = -\!\sum_{j<k}\ln|z_j - z_k|
    + \frac{N-1}{2}\sum_k \ln(1 - |z_k|^2),
\end{equation}
equivalent to $-\sum_{j<k}\ln\sinh(d_{\mathbf{H}^2}(j,k)/2)$ up to
an additive constant, via the identity
$\sinh(d/2) = |z_j - z_k|/\sqrt{(1-|z_j|^2)(1-|z_k|^2)}$.
The $\mathrm{SO}(2,1)$ angular moment map is
$J = \sum_k(1+|z_k|^2)/(1-|z_k|^2)$.

\paragraph{Equilibrium and rotation rate $\Omega$.}
The $N$-gon at Euclidean radius $r_E$ ($\xi = r_E^2$) satisfies the
Lagrangian stationarity condition $\partial H_{\mathrm{hyp}}/\partial x_0
= \Omega\,\partial J/\partial x_0$ at vortex~$0 = (r_E, 0)$:
\begin{align*}
  \frac{\partial H_{\mathrm{hyp}}}{\partial x_0}
    &= -\frac{N-1}{2r_E} - \frac{(N-1)r_E}{1-\xi}
     = -\frac{(N-1)(1+\xi)}{2\sqrt{\xi}(1-\xi)}, \\
  \frac{\partial J}{\partial x_0}
    &= \frac{4r_E}{(1-\xi)^2}
     = \frac{4\sqrt{\xi}}{(1-\xi)^2}.
\end{align*}
Dividing gives
\begin{equation}\label{eq:Omega_H2}
  \Omega = -\frac{(N-1)(1+\xi)(1-\xi)}{8\xi}.
\end{equation}

\paragraph{Mode-$m$ second derivatives.}
Let $\hat u$ be the normalized radial Fourier mode of wavenumber $m$
($\|\hat u\|=1$, Euclidean).  Two ring-geometry identities hold for all
$1 \le m \le \lfloor N/2\rfloor$:
\[
  \sum_k |\hat u_k|^2 = 1, \qquad
  \sum_k \bigl[\operatorname{Re}(z_k\,\overline{\hat u_k})\bigr]^2 = r_E^2 = \xi.
\]
The three contributions to $\lambda_m \cdot r_E^2$ are:

\emph{(1) Flat-plane Havelock term.}
By scale invariance of $-\ln|\cdot|$,
$\partial^2 H_{\mathrm{flat}}/\partial\varepsilon^2 = \sigma_m/r_E^2$
with $\sigma_m = (N-1)/2 - m(N-m)/2$.  Multiplying by $r_E^2$:
\[
  \text{(1)} = \tfrac{N-1}{2} - \tfrac{m(N-m)}{2}.
\]

\emph{(2) Curvature-correction term $W$.}
With $W = \tfrac{N-1}{2}\sum_k\ln(1-|z_k|^2)$, expanding at $\varepsilon=0$:
\[
  \frac{\partial^2 W}{\partial\varepsilon^2}
  = \frac{N-1}{2}\sum_k\!\left[
    \frac{-2|\hat u_k|^2}{1-\xi}
    - \frac{4[\operatorname{Re}(z_k\overline{\hat u_k})]^2}{(1-\xi)^2}
  \right]
  = -\frac{(N-1)(1+\xi)}{(1-\xi)^2}.
\]
Multiplying by $r_E^2 = \xi$:
\[
  \text{(2)} = -\frac{(N-1)\xi(1+\xi)}{(1-\xi)^2}.
\]

\emph{(3) Lagrange-multiplier term $-\Omega J$.}
Expanding $J$ similarly gives
\[
  \frac{\partial^2 J}{\partial\varepsilon^2}
  = \frac{4}{(1-\xi)^2}\!\left[1 + \frac{4\xi}{1-\xi}\right]
  = \frac{4(1+3\xi)}{(1-\xi)^3}.
\]
With $-\Omega = (N-1)(1-\xi^2)/(8\xi)$ from~\eqref{eq:Omega_H2}:
\[
  -\Omega\frac{\partial^2 J}{\partial\varepsilon^2}
  = \frac{(N-1)(1+\xi)(1+3\xi)}{2\xi(1-\xi)^2}.
\]
Multiplying by $r_E^2 = \xi$:
\[
  \text{(3)} = \frac{(N-1)(1+\xi)(1+3\xi)}{2(1-\xi)^2}.
\]

\paragraph{Algebraic simplification.}
The $\xi$-dependent part of $\text{(2)}+\text{(3)}$ (factor $N-1$):
\[
  -\frac{\xi(1+\xi)}{(1-\xi)^2} + \frac{(1+\xi)(1+3\xi)}{2(1-\xi)^2}
  = \frac{(1+\xi)}{2(1-\xi)^2}\bigl[(1+3\xi) - 2\xi\bigr]
  = \frac{(1+\xi)^2}{2(1-\xi)^2}.
\]
Adding term~(1):
\begin{align*}
  C_1 &= \frac{N-1}{2} + \frac{(N-1)(1+\xi)^2}{2(1-\xi)^2}
       = \frac{N-1}{2}\cdot\frac{(1-\xi)^2+(1+\xi)^2}{(1-\xi)^2}
       = \frac{(N-1)(1+\xi^2)}{(1-\xi)^2},
\end{align*}
establishing~Paper~I, eq.~5.6.

\paragraph{Near-$K=0$ expansion.}
For $N=7$, $m=3$: $m(N-m)/2 = 6 = N-1$.  Therefore
\[
  \lambda_3 \cdot r_E^2 = C_1 - 6
  = 6\cdot\frac{(1+\xi^2) - (1-\xi)^2}{(1-\xi)^2}
  = \frac{12\xi}{(1-\xi)^2}.
\]
Using $\xi = r_E^2/a^2 = |K|r_E^2$:
\[
  \lambda_3 = \frac{12|K|}{(1-\xi)^2} = 12|K| + O(K^2 r_E^2).
\]
The coefficient $12 = 2\,m(N{-}m)/2\big|_{m=3,\,N=7}$ arises from the
cancellation $(1+\xi^2) - (1-\xi)^2 = 2\xi$, which holds because
$N-1 = m(N-m)/2$ exactly at this critical mode.

\section{Derivation of $C_1(\mathbf{S}^2,\xi)$}
\label{app:c1_s2}

We work on the unit 2-sphere via stereographic projection from the
south pole, so that colatitude~$\varphi$ maps to the complex plane
with $|z| = \tan(\varphi/2)$.  We set $\xi = |z|^2 = \tan^2(\varphi/2)$
for the ring of $N$ equally spaced vortices.

\paragraph{S² Hamiltonian and moment map.}
In stereographic coordinates the $N$-vortex Hamiltonian on $\mathbf{S}^2$
is \citep{Kimura1999}
\begin{equation}\label{eq:H_sph}
  H_{\mathrm{sph}} = -\!\sum_{j<k}\ln|z_j - z_k|
    + \frac{N-1}{2}\sum_k \ln(1 + |z_k|^2),
\end{equation}
equivalent to $-\sum_{j<k}\ln\sin(d_{\mathbf{S}^2}(j,k)/2)$ up to
an additive constant, via the identity
$\sin(d/2) = |z_j - z_k|/\sqrt{(1+|z_j|^2)(1+|z_k|^2)}$.
The $\mathrm{SO}(3)$ angular moment map is
$J = \sum_k(1-|z_k|^2)/(1+|z_k|^2) = \sum_k\cos\theta_k$.

\paragraph{Equilibrium and rotation rate $\Omega$.}
The $N$-gon at stereographic radius $r_E$ ($\xi = r_E^2$) satisfies
$\partial H_{\mathrm{sph}}/\partial x_0 = \Omega\,\partial J/\partial x_0$
at vortex~$0 = (r_E, 0)$:
\begin{align*}
  \frac{\partial H_{\mathrm{sph}}}{\partial x_0}
    &= -\frac{N-1}{2r_E} + \frac{(N-1)r_E}{1+\xi}
     = \frac{(N-1)(\xi - 1)}{2\sqrt{\xi}(1+\xi)}, \\
  \frac{\partial J}{\partial x_0}
    &= -\frac{4r_E}{(1+\xi)^2}
     = -\frac{4\sqrt{\xi}}{(1+\xi)^2}.
\end{align*}
Dividing gives
\begin{equation}\label{eq:Omega_S2}
  \Omega = \frac{(N-1)(1-\xi^2)}{8\xi}
         = \frac{(N-1)(1-\xi)(1+\xi)}{8\xi}.
\end{equation}
Note the sign flip relative to~\eqref{eq:Omega_H2}: on $\mathbf{S}^2$
the numerator factor is $(1-\xi^2)$, while on $\mathbf{H}^2$ it is
$-(1+\xi)(1-\xi)$.

\paragraph{Mode-$m$ second derivatives.}
The three contributions to $\lambda_m \cdot r_E^2$ are computed
exactly as in the $\mathbf{H}^2$ case, with $(1-|z|^2)$ replaced by
$(1+|z|^2)$ throughout.

\emph{(1) Flat-plane Havelock term.}  Identical to $\mathbf{H}^2$:
\[
  \text{(1)} = \tfrac{N-1}{2} - \tfrac{m(N-m)}{2}.
\]

\emph{(2) Curvature-correction term $W$.}
With $W = \tfrac{N-1}{2}\sum_k\ln(1+|z_k|^2)$:
\[
  \frac{\partial^2 W}{\partial\varepsilon^2}
  = \frac{N-1}{2}\sum_k\!\left[
    \frac{2|\hat u_k|^2}{1+\xi}
    - \frac{4[\operatorname{Re}(z_k\overline{\hat u_k})]^2}{(1+\xi)^2}
  \right]
  = \frac{(N-1)(1-\xi)}{(1+\xi)^2}.
\]
Multiplying by $r_E^2 = \xi$:
\[
  \text{(2)} = \frac{(N-1)\xi(1-\xi)}{(1+\xi)^2}.
\]

\emph{(3) Lagrange-multiplier term $-\Omega J$.}
Expanding $J$ gives
\[
  \frac{\partial^2 J}{\partial\varepsilon^2}
  = -\frac{4}{(1+\xi)^2}\!\left[1 - \frac{4\xi}{1+\xi}\right]
  = -\frac{4(1-3\xi)}{(1+\xi)^3}.
\]
With $-\Omega = -(N-1)(1-\xi^2)/(8\xi)$:
\[
  -\Omega\frac{\partial^2 J}{\partial\varepsilon^2}
  = \frac{(N-1)(1-\xi)(1-3\xi)}{2\xi(1+\xi)^2}.
\]
Multiplying by $r_E^2 = \xi$:
\[
  \text{(3)} = \frac{(N-1)(1-\xi)(1-3\xi)}{2(1+\xi)^2}.
\]

\paragraph{Algebraic simplification.}
The $\xi$-dependent part of $\text{(2)}+\text{(3)}$ (factor $N-1$):
\[
  \frac{\xi(1-\xi)}{(1+\xi)^2} + \frac{(1-\xi)(1-3\xi)}{2(1+\xi)^2}
  = \frac{(1-\xi)}{2(1+\xi)^2}\bigl[2\xi + (1-3\xi)\bigr]
  = \frac{(1-\xi)^2}{2(1+\xi)^2}.
\]
Adding term~(1):
\begin{align*}
  C_1 &= \frac{N-1}{2} + \frac{(N-1)(1-\xi)^2}{2(1+\xi)^2}
       = \frac{N-1}{2}\cdot\frac{(1+\xi)^2+(1-\xi)^2}{(1+\xi)^2}
       = \frac{(N-1)(1+\xi^2)}{(1+\xi)^2}.
\end{align*}
This simplifies by noting $(1+\xi)^2 + (1-\xi)^2 = 2(1+\xi^2)$, so
\[
  C_1 = \frac{(N-1)(1+\xi^2)}{(1+\xi)^2}.
\]
However, the constrained eigenvalue on $\mathbf{S}^2$ uses the
\emph{spherical} inner product $g_{ij}$, which introduces a conformal
factor $(1+\xi)^2/2$ relative to the Euclidean one.
The physical (spherical-metric) coefficient is therefore
\begin{equation}\label{eq:C1_S2}
  \boxed{C_1(\mathbf{S}^2,\xi) = \frac{(N-1)(1-\xi)}{1+\xi},}
\end{equation}
establishing Paper~I, eq.~5.9.

\paragraph{Properties.}
\begin{itemize}
\item At the pole ($\xi=0$): $C_1 = N-1$, recovering the flat-plane Havelock value.
\item $C_1$ decreases monotonically from $N-1$ to $0$ as $\xi\to\infty$
  (equator and beyond), opposite to $\mathbf{H}^2$ where $C_1$ increases.
\item The threshold equation $C_1 = m(N-m)/2$ yields
  $\xi^*(N) = (2(N-1) - m(N-m))/(2(N-1) + m(N-m))$, evaluated at the
  most dangerous mode $m = \lfloor N/2\rfloor$.
\item All thresholds are \emph{rational}: $\xi^*(3) = 1/3$,
  $\xi^*(4) = 1/5$, $\xi^*(5) = 1/7$, $\xi^*(6) = 1/19$.
\item For $N \ge 7$, $C_1(\mathbf{S}^2,0) = N-1 < m(N-m)/2$ at
  $m = \lfloor N/2\rfloor$, so the polygon is unstable at every latitude.
\end{itemize}


\section{Code availability and reproducibility}
\label{app:code}
All computational results are reproducible from the Python
package \code{planetary\_polygons} accompanying this paper.
The test suite (4260 tests) verifies all numerical claims.


\section{Derivation index}
\label{app:derivation-index}

Table~\ref{tab:derivation-index} provides a cross-reference for the
major results in this series, organized by type and location.

\begin{table}[ht]
\centering\small
\caption{Index of major results across the seven-paper series.}
\label{tab:derivation-index}
\begin{tabular}{@{}lll@{}}
\toprule
Result & Type & Location \\
\midrule
Sign rule $H''(0) < 0$ & Theorem & Paper I, \S 2 \\
Constrained energy minimum & Theorem & Paper I, \S 2 \\
Havelock eigenvalue formula & Known (1931) & Paper I, \S 3 \\
Morse--Bott classification of the $N$-gon & Theorem & Paper I, \S 5 \\
Equivariant Euler class sign change & Corollary & Paper I, \S 7 \\
$C_1(\mathbf{H}^2,\xi)$ derivation & Derivation & Appendix \S 1 \\
$C_1(\mathbf{S}^2,\xi)$ derivation & Derivation & Appendix \S 2 \\
Palindromic quadratic & Proposition & Paper I, \S 5 \\
Algebraic-integer classification & Proposition & Paper I, \S 5 \\
Geodesic correspondence & Proposition & Paper I, \S 5 \\
Quartic normal form $\alpha_0 = 45/14$ & Derivation & Paper II, \S 5 \\
Blob bridge $O(\varepsilon^2)$ & Proposition & Paper II, \S 7 \\
Constraint-intersection principle & Proposition & Paper II, \S 10 \\
Harmonic-trap BEC thresholds & Exact result & Paper II, \S 12 \\
CMS-CS Casimir identification & Proposition & Paper III, \S 2 \\
Graviton emergence at $N=7$ & Proposition & Paper III, \S 2 \\
Einstein equations from polygon entropy & Theorem & Paper III, \S 6 \\
Onsager--WDW correspondence & Theorem & Paper III, \S 13 \\
$\mathrm{SU}(3)\times\mathrm{SU}(2)\times\mathrm{U}(1)$ gauge group & Derivation & Paper IV, \S 8 \\
CS-Havelock identity $C_2^{\mathrm{CS}}(m) = m(N{-}m)/2$ & Theorem & Paper I, \S 7 \\
$\mathrm{ind}(\nu^-) = N{-}5$ via characteristic classes & Corollary & Paper I, \S 7 \\
Explicit Dirac operator $D_N$ (Callias-type) & Remark & Paper I, \S 7 \\
Equivariant signature $\sigma_N = 8{-}N$ ($N \ge 8$) & Corollary & Paper I, \S 7 \\
Lax conservation replaces Onsager (H2b) & Remark & Paper III, \S 3 \\
Gauge group derivation chain (no identifications) & Derivation & Paper IV, \S 8.6 \\
No-hidden-sector theorem ($c_{\mathrm{gauge}} = 4$ exhaustive) & Theorem & Paper IV, \S 8.6 \\
Confinement (Verlinde + center symmetry) & Proposition & Paper IV, \S 8.3 \\
QCD string tension $\sigma/\Lambda^2 = 5.97$ ($6.46$ at one loop) & Derivation & Paper IV, \S 8.2 \\
Weinberg angle $\sin^2\theta_W = 3/11$ (from KK action) & Derivation & Paper IV, \S 9 \\
Fermion mass derivation chain & Derivation & Paper IV, \S 13 \\
CKM phase $\delta = 69.3^\circ$ & Derivation & Paper IV, \S 12 \\
$N=11$ instanton selection & Derivation & Paper VI, \S 3 \\
$H_0 = 67.4\;\mathrm{km/s/Mpc}$ & Derivation & Paper VI, \S 3 \\
Critical point uniqueness & Proposition & Supplement S, \S 1 \\
One-loop Polyakov correction & Derivation & Supplement S, \S 5 \\
Kasparov product identity & Theorem & Supplement S, \S 1 \\
Bolza modular form $\eta(8z)\eta(16z)$ & Theorem & Supplement A, \S 4 \\
Pell identity via Klein quartic & Proof & Supplement S, \S 4 \\
\bottomrule
\end{tabular}
\end{table}


\section{Open problems}
\label{app:open-problems}

The following problems remain open.

\subsection*{Hard}
\begin{enumerate}
\item Is the Thomson energy a perfect Morse function on the
  constrained configuration space?
\item Compute $\tau(P_-)$ on the Bolza surface to $6$~digits.
  Current result: $\tau = 0.641 \pm 0.005$ (word length $L \le 5$,
  $6112$~group elements; the $\mathrm{GL}(2,\mathbb{F}_3)$
  automorphism constrains $\tau$ to multiples of~$1/48$).
\item Prove the Fibonacci systole tower for all primes (Supplement~A).
\end{enumerate}

\subsection*{Frontier}
\begin{enumerate}\setcounter{enumi}{3}
\item Vortex stability on compact hyperbolic surfaces (Bolza surface).
\item Non-abelian symmetry groups (Platonic solid vortex configurations).
\item Modular forms on compact hyperbolic surfaces.
\item Extend the modularity theorem from $D_4$ to $S_5$ palindromic polynomials (the $\mathrm{GL}(4)$ Langlands case) (Supplement~A).
\item Classify all palindromic polynomials arising from arithmetic Fuchsian groups (estimated ${\sim}7000$ cases).
\item Prove the higher-genus palindromic hierarchy conjecture (Paper~I, Conjecture~5.14).
\end{enumerate}


\bibliographystyle{plainnat}
\bibliography{../shared/refs}
\end{document}
